\documentclass[book.tex]{subfiles}

\begin{document}

\chapter{Phonon Gas}

\label{chap:phonons}
\noindent\rule{11cm}{0.4pt} \\
\textbf{Prerequisites:}  \autoref{chap:ising_model}  \\
\textbf{Difficulty Level:} ** \\
\noindent\rule{11cm}{0.4pt}
\vspace{5mm}

Analogous to the relationship of photons to electromagnetic or light energy, phonons are a definite discrete unit of quantum vibrational mechanical energy. Phonons and electrons transport are the main mechanism of energy diffusion in crystals. Thus for non-conducting crystals, heat transport is mostly carried by phonons transport, and so a proper description of the thermodynamic of phonons is necessary to understand the origins of the thermal properties of solids.

In this chapter, we use statistical mechanics to derive the thermodynamic of phonon gases, which are the normal modes of low temperature solids. At low temperatures, atoms in a crystal remain close to their equilibrium position, so a simple model can be developed to describe the energy states of this particles. At the end of this chapter we will derive the heat capacity of crystals due to the contribution of phonons considering two different approaches: the \textit{Einstein Model} which consider all phonons vibrating at the same frequency (valid for optical phonons); and the \textit{Debye Model} which consider a linear dispersion of frequencies (valid for acoustic phonons at low temperatures). 

\section{Crystalline Solid}

\begin{marginfigure}
% \includegraphics[width=\linewidth]{figures/Physics/statistical_mechanics/phonon_gas/Fig4.png}
\includegraphics[width=\linewidth]{figures/Physics/statistical_mechanics/phonon_gas/PhononGas4.png}
\caption{Crystal lattice structures}
\label{fig4}
\end{marginfigure}

The atoms in a crystal are arranged in a regular array of points in space with a variety of possible crystalline lattices (Fig.~\ref{fig4}). At zero temperature, the atomic coordinates are uniquely locked into spatial positions that minimize the energy. At finite temperature, the atoms fluctuate about the energy-minimum positions, leading to lattice vibrations that govern the thermodynamic behavior (Fig.~\ref{fig5}).

\begin{figure}
% \includegraphics[width=\linewidth]{figures/Physics/statistical_mechanics/phonon_gas/Fig5.png}
\includegraphics[width=\linewidth]{figures/Physics/statistical_mechanics/phonon_gas/atoms _in_crystel5.png}

\caption{Schematic showing the atoms in a crystal in their locked, energy-minimum positions at $T = 0$ and the lattice vibrations that occur at finite temperature}
\label{fig5}
\end{figure}

\section{Thermodynamic contribution of lattice vibrations}

Consider N atoms with positions 
\begin{align}
\{\vec{r}\} = \vec{r}_1, \vec{r}_2,\dotsc,\vec{r}_N
\end{align}
in a crystalline lattice. Define the potential energy $V\left(\{\vec{r}\}\right)$ that describes the energy for a given system configuration  $\{\vec{r}\}$. As before, we will simplify notation to $V(\vec{r})$ when the meaning is clear from context. A minimum-energy configuration 
\begin{align}
\{\vec{r}^{(0)}\} &= \vec{r}^{(0)}_1, \vec{r}^{(0)}_2,\dotsc,\vec{r}^{(0)}_N
\end{align}
satisfies the condition 
\begin{align}
\nabla_{\vec{r}_i} V\left(\vec{r}^{(0)}\right) = \nabla_{\vec{r}_i} V|_0 = \vec{0}
\end{align} 
for $i = 1,2,...,N$. The atomic positions $\{\vec{r}^{(0)}\}$ define the regular crystalline lattice.

Expand the potential energy about $\{\vec{r}\} = \{\vec{r}^{(0)}\}$, such that:
\begin{align}
V\left(\vec{r}\right) \approx \  &V\left(\vec{r}^{(0)}\right) + \sum_{i=1}^N \left(\nabla_{\vec{r}_i} V|_0\right)\cdot \left(\vec{r}_i - \vec{r}^{(0)}_i \right) + \notag\\ &\frac{1}{2} \sum_{i=1}^N \sum_{j=1}^N \left(\nabla_{\vec{r}_i} \nabla_{\vec{r}_j} V|_0\right) \left(\vec{r}_i - \vec{r}^{(0)}_i \right) \left(\vec{r}_j - \vec{r}^{(0)}_j \right) + \dotsc \notag
\end{align}

The linear term is zero (by definition), thus the energy is:
\begin{align}\label{Eq17}
V \approx V_0 + \frac{1}{2} \sum_{i=1}^N \sum_{j=1}^N \sum_{\alpha,\gamma = x,y,z} s_{i\alpha}s_{j\gamma}K_{i\alpha j \gamma}
\end{align}
where 
\begin{align}
s_{i\alpha} &= \hat{e}_\alpha \cdot \left(\vec{r}_i - \vec{r}^{(0)}_i \right) \\ 
K_{i\alpha j \gamma} &= \nabla_{\vec{r}_i} \nabla_{\vec{r}_j} V|_0
\end{align}

The matrix $K_{i\alpha j \gamma}$ can be diagonalized into normal modes (eigenvectors) with effective elastic constants (eigenvalues). Since there are $3N-6\approx 3N$ degrees of freedom, there are $3N$ normal modes.

The potential energy is written as:
\begin{align}\label{Eq18}
V = V_0 + \frac{1}{2}\sum_{l=1}^{3N} \tilde{K}_l \xi_l^2
\end{align}

\noindent where $\xi_l$ is the magnitude of the $l$th normal mode.

The total energy of the system is then written as:
\begin{align}\label{Eq19}
E = \sum_{l=1}^{3N}\frac{\tilde{p}_l^2}{2\tilde{m}_l} + \frac{1}{2}\sum_{l=1}^{3N}\tilde{K}_l \xi_l^2 + V_0
\end{align}

\noindent where $\tilde{p}_l$ and $\tilde{m}_l$ are the effective momentum and mass of the $l$th normal mode, respectively.

The total energy $E$ is decomposed into normal modes with an individual-mode energy in the form of a harmonic oscillator. These normal modes are called \textit{phonons}. Phonons act as quasi-particles which means they are distinguishable and independent, \textit{i.e.} they don't interact between each other.

The Hamiltonian (sum of kinetic and potential energy) of the $l$th
phonon is given by 
\begin{align}
H_l &= \frac{\tilde{p}_l^2}{2\tilde{m}_l} + \frac{1}{2}K_l\xi^2_l
\end{align}
This harmonic oscillator energy results in the quantized energy:
\begin{align}\label{Eq20}
E_l = \left(j_l + \frac{1}{2}\right)\hslash\omega_l
\end{align}

\noindent where $j_l = 0,1,2...$, and the phonon frequency is $\omega_l = \sqrt{\frac{\tilde{K}_l}{\tilde{m}_l}}$.

The canonical partition function $Q$ is given by:
\begin{align}\label{Eq21}
Q(T,V,N) &= \sum_{j_1 = 0}^\infty \sum_{j_2 = 0}^\infty ... \sum_{j_{3N} = 0}^\infty e^{\left[ -\beta V_0 - \beta \sum_{l=1}^{3N} \left(j_l + \frac{1}{2}\right)\hslash\omega_l \right]} \notag\\
&= e^{ -\beta V_0 } \prod_{l=1}^{3N} \sum_{j_l = 0}^\infty e^{-\beta\left(j_l + \frac{1}{2}\right)\hslash\omega_l} \notag\\
&=e^{ -\beta V_0 } \prod_{l=1}^{3N} \frac{e^{-\beta\frac{1}{2}\hslash\omega_l}}{1 - e^{-\beta\hslash\omega_l}}
\end{align}

\noindent where we use the mathematical property \begin{align}
\sum_{n=0}^\infty z^n &= \frac{1}{1 - z}
\end{align}

This gives the Helmholtz free energy:
\begin{align}
F = -k_BT\log Q = V_0 + k_BT \sum_{l=1}^{3N}\left[\log\left(1-e^{-\beta\hslash\omega_l}\right) + \frac{1}{2}\beta\hslash\omega_l \right] \notag
\end{align}

The average energy is given by:
\begin{align}\label{Eq.E_Phonons}
\langle E \rangle = -\frac{\partial \log Q}{\partial \beta} = V_0 + \sum_{l=1}^{3N} \left( \frac{\hslash\omega_le^{-\beta\hslash\omega_l}}{1 - e^{-\beta\hslash\omega_l}} + \frac{1}{2}\hslash\omega_l \right)
\end{align}

The next step is to analyze two competing models for the phonon frequencies $\omega_l$.

\section{Einstein model}
In the Einstein model, we assume there is a single characteristic
frequency of the crystal, defined as the Einstein frequency $\omega_E$. This model is a good approximation for optical phonons.

The average energy for this model is given by:
\begin{align}
\langle E \rangle = V_0 + \frac{3N\hslash\omega_E}{2} + \frac{3N\hslash\omega_Ee^{-\beta\hslash\omega_l}}{1-e^{-\beta\hslash\omega_E}} \notag	
\end{align}

For this model, the heat capacity is given by:
\begin{align}
C_V = \left( \frac{\partial \langle E \rangle}{\partial T} \right)_{V,N} =  \frac{3Nk_B \left( \frac{\hslash\omega_E}{k_BT} \right) ^2 e^{-\frac{\hslash\omega_E}{k_BT}}}{\left(1-e^{\frac{-\hslash\omega_E}{k_BT}}\right)^2} = \frac{3Nk_B \left( \frac{\theta_E}{T} \right) ^2 e^{-\frac{\theta_E}{T}}}{\left(1-e^{\frac{\theta_E}{T}}\right)^2} \notag
\end{align}

\noindent where we define 
\begin{align}
\theta_E = \frac{\hslash\omega_E}{k_B}
\end{align}

In the limit $T \rightarrow \infty$ the heat capacity 
\begin{align}
C_V \rightarrow 3Nk_B
\end{align}
This is known as the Dulong-Petit Law. As we learned before, the equipartition theorem states that the energy receives $k_BT$ per thermally active degree of freedom for a harmonic oscillator (quantized energy is linear in the quantum index). In the limit $T \rightarrow 0$ the heat capacity 
\begin{align}
C_V \rightarrow 0
\end{align}
The heat capacity approaches zero exponentially in the small-$T$ limit (Fig.~\ref{fig6}).

\begin{marginfigure}
\includegraphics[width=\linewidth]{figures/Physics/statistical_mechanics/phonon_gas/Fig6.png}
\caption{Heat capacity of a crystal predicted by the Einstein model of lattice vibration}
\label{fig6}
\end{marginfigure}

\section{Debye model}
The Debye model treats the solid as an elastic material. This approximation is usually accurate to describe acoustic phonons at low temperatures.

Vibrational modes in an elastic solid correspond to sound waves, thus the frequencies satisfy 
\begin{align}
\omega=ck
\end{align} 
where $c$ is the speed of sound in the solid and \begin{align}
k = m\pi/L
\end{align}
where $m=1,2,\dotsc$.

Thus, for considering a long particle chain, we convert the sum over normal modes into a integral over the frequencies $\omega$ using:
\begin{align}
\sum_{l=1}^{3N} \left( \frac{1}{e^{\beta\hslash\omega_l} - 1} + \frac{1}{2} \right)\hslash\omega_l &=\sum_{m_1}\sum_{m_2}\sum_{m_3} \left( \frac{1}{e^{\beta\hslash\omega_l} - 1} + \frac{1}{2} \right)\hslash\omega_l \notag\\
&= \iiint \left[\left( \langle n \rangle + \frac{1}{2} \right)\hslash\omega_l\right] dm_1dm_2dm_3 \notag\\
&=\iiint_0^{k_c} \left(\frac{L}{\pi}\right)^3 \left[\left( \langle n \rangle + \frac{1}{2} \right)\hslash\omega_l\right] dk_1dk_2dk_3 \notag\\
&= \int_0^{k_c} 4\pi\left(\frac{L}{\pi}\right)^3 \left( \langle n \rangle + \frac{1}{2} \right)\hslash\omega_l k^2 dk \notag\\
&= \int_0^{\omega_D}D(\omega) \left( \langle n \rangle + \frac{1}{2} \right) \hslash\omega_l d\omega\notag
\end{align}

where 
\begin{align}
D(\omega) = \frac{4L^3\omega^2}{\pi^2c^3}
\end{align}
is the density of states of acoustic phonons , and 
\begin{align} 
\langle n \rangle = \frac{1}{e^{\beta\hslash\omega_l} - 1}
\end{align}
is the Bose-Einstein Distribution. In this equation $k_c$ is a cutoff wavemode (to be determined), and $\omega_D = Ck_c$ is called the Debye frequency.

A complete conversion will include one longitudinal mode with $c_l$ and two transverse modes with $c_t$ . This gives:
\begin{align}\label{Eq22}
\sum_l \left( \frac{1}{e^{\beta\hslash\omega_l} - 1} + \frac{1}{2} \right)\hslash\omega_l &= \int_0^{\omega_D} \frac{4L^3\omega^2}{\pi^2}\left(\frac{1}{c_l^3} + \frac{2}{c_t^3} \right) \left( \langle n \rangle + \frac{1}{2} \right) \hslash\omega_l d\omega \notag\\
&=\int_0^{\omega_D} D(\omega) \left( \langle n \rangle + \frac{1}{2} \right) \hslash\omega_l d\omega
\end{align}

\noindent where we defined a new form for the density of states given by 
\begin{align} 
D(\omega) = \frac{4L^3\omega^2}{\pi^2}\left(\frac{1}{c_l^3} + \frac{2}{c_t^3} \right)
\end{align}

To find $\omega_D$ we must enforce that 
\begin{align}
\sum_l 1 &= 3N,
\end{align}
Thus we have:
\begin{align}
\sum_l 1 &= \int_0^{\omega_D}D(\omega) \left( \langle n \rangle + \frac{1}{2} \right) d\omega \notag\\
&= \left(\frac{L}{\pi}\right)^3 \left(\frac{1}{c_l^3} + \frac{2}{c_t^3} \right) \frac{\omega_D^3}{3} = 3N \notag\\&\Rightarrow \omega_D = \frac{\pi}{L}\left[\frac{9N}{\left(\frac{1}{c_l^3} + \frac{2}{c_t^3} \right)}\right]^{1/3} \notag
\end{align}

For convenience, use $\omega_D$ as a parameter, thus we can define the density of states as:
\begin{align}\label{Eq23}
D(\omega) = \frac{9N\omega^2}{\omega_D}
\end{align}

\begin{marginfigure}
\includegraphics[width=\linewidth]{figures/Physics/statistical_mechanics/phonon_gas/Fig7.png}
\caption{Heat capacity of a crystal predicted by the Einstein model and Debye model of lattice vibration}
\label{fig7}
\end{marginfigure}

Define the Debye temperature 
\begin{align}
\theta_D = \frac{\hslash\omega_D}{k_B}
\end{align}
which defines the temperature scale for vibrational fluctuations. To test this model, we find the heat capacity:
\begin{align}
C_V &= -k_B\beta^2\left(\frac{\partial\langle E\rangle}{\partial\beta}\right) \\
&= 9Nk_B\frac{T^3}{\theta_D^3}\int_0^{\theta_D/T}\frac{x^4e^{-x}}{\left(1-e^{-x}\right)^2} dx 
\end{align}

In the limit $T \rightarrow \infty $ the heat capacity approaches:
\begin{align}
C_V \rightarrow 9Nk_B\frac{T^3}{\theta_D^3}\int_0^{\theta_D/T}dxx^2 = 3Nk_B \notag
\end{align}

\noindent which is expected (Dulong-Petit Law), since at high temperatures all phonon modes become activated.

In the limit $T \rightarrow 0 $ the heat capacity scales as:
\begin{align}
C_V \rightarrow 9Nk_B\frac{T^3}{\theta_D^3}\int_0^{\infty}\frac{x^4e^{-x}}{\left(1-e^{-x}\right)^2} dx \sim Nk_B\frac{T^3}{\theta_D^3}
\end{align}

To obtain a qualitative picture of the Debye $T^3$ law, we suppose that all phonon modes of wavevector less than $K_T$ have the classical thermal energy
$k_BT$ and that modes between $K_T$ and the Debye cutoff $K_D$ are not excited at all. Of the $3N$ possible modes, the fraction excited is 
\begin{align} 
(K_T/K_D)^3 &=(T/\theta)^3
\end{align}
because this is the ratio of the volume of the inner sphere to the outer sphere. The energy is 
\begin{align} 
U \approx k_BT \cdot 3N(T/\theta)^3
\end{align}
and the heat capacity is 
\begin{align} 
C_v &=\frac{\partial U}{\partial T} = 12Nk_B(T/\theta)^3
\end{align}


The heat capacity predicted by the Einstein and the Debye models are very similar. However, the low temperature of the heat capacity of non-conducting solids matches the Debye model, \textit{i.e.} $C_V \sim NT^3$ (Fig.~\ref{fig7}). This is because at low temperatures only acoustic phonons are activated, which is better represented in the Debye model.

\end{document}

